\pdfminorversion=7
\documentclass[11pt,a4paper]{article}

\usepackage[T1]{fontenc}
\usepackage[utf8]{inputenc}
\usepackage{lmodern}
\usepackage{microtype}
\usepackage[margin=1in]{geometry}
\usepackage{mathtools,amssymb,bm}
\usepackage{graphicx}
\usepackage{enumitem}
\usepackage[authoryear,round]{natbib}
\usepackage[hidelinks]{hyperref}

\hypersetup{
  pdftitle={Refractive Gravity: An Ontological Picture and the Foundations of Mathematical Formalism},
  pdfauthor={Giorgi Lotishvili}
}

\graphicspath{{figures/}}
\setcounter{tocdepth}{2}
\setlist{itemsep=0.35em,topsep=0.5em}

\newcommand{\RefG}{\mathrm{RefG}}

\title{\textbf{Refractive Gravity: An Ontological Picture and the Foundations of Mathematical Formalism}}
\author{Giorgi Lotishvili\\
\small Associate Professor, East European University\\
\small Tbilisi, Georgia\\
\small \href{mailto:g_lotishvili@eeu.edu.ge}{\texttt{g\_lotishvili@eeu.edu.ge}}}
\date{2026}

\begin{document}

\maketitle

\begin{center}
\textbf{Zenodo DOI:} \href{https://doi.org/10.5281/zenodo.20633284}{10.5281/zenodo.20633284}
\end{center}

\begin{abstract}
This article presents an ontological account of \emph{Refractive Gravity} (RefG), bringing several scientific works and lines of inquiry together within a single physical vision. This is not a static theory, but a research program developing step by step. The principal aim of the text is to offer the reader an intuitive picture of the fundamental organization of the Universe; it is therefore deliberately not overloaded with formulas. The current rigorous formalism and completed calculations are available separately in the public Zenodo research archive.

For the current status of the ideas and results discussed in this article, see the final section, ``Current Status of the Research Program.''
\end{abstract}

\tableofcontents
\clearpage

\section*{Introduction}
\addcontentsline{toc}{section}{Introduction}

Modern physics resembles a complex puzzle: its individual pieces have been calculated to high precision, tested experimentally, and refined mathematically, yet they still do not fit together into a unified physical picture. The Standard Model of particle physics and the geometric description of gravity have not yet been fully brought within a common language; in the interiors of black holes, classical geometry leads to singularities; in galactic dynamics, rotation curves and the regularities of modified Newtonian dynamics (MOND) are explained within standard cosmology by dark-matter halos, while the accelerated expansion of the Universe is described in the language of dark energy. This broader set of questions also includes quantum randomness, the nature of time, the hierarchy of particle masses, the fundamental constants, and the relative weakness of gravity. Each has been studied in detail, but their possible common ontological foundation remains unclear.

Throughout the history of physics, these problems have elicited many unifying visions. For Albert Einstein, the path to unification lay in a deeper geometry of spacetime. John Archibald Wheeler sought a fundamental picture of the Universe at the boundary between geometry and information through his principle of ``It from Bit.'' Andrei Sakharov interpreted gravity as an effect emerging from quantum fluctuations of the vacuum; Ted Jacobson derived the Einstein equation from the thermodynamic regularities of local horizons; and Erik Verlinde viewed gravity as an entropic, holographic phenomenon. Roger Penrose, in turn, treats the global structure of spacetime, cosmology, and quantum physics as a unified problem of twistor and cyclic geometry. These approaches differ in their mechanisms, but they ask a common question: can different phenomena be described as arising from a single foundation?

This article organizes ideas developed over many years into a unified logical structure and presents the ontological framework of \textbf{Refractive Gravity (RefG)}. The standard language of modern physics generally begins with a ready-made apparatus of measurement: coordinates, a time parameter, and fields are taken as given starting points. RefG begins at the threshold that precedes this apparatus. Before we can speak of space, time, or mass, nature must contain a physically distinguishable, persistent difference. What physical condition gives these concepts the form of measurable quantities?

In 1905, Einstein abandoned absolute time and defined simultaneity operationally through clocks and light signals \citep{einstein1905}. RefG takes this operational question one level deeper: what physical condition creates the clock and signal on which those operations themselves depend?

To answer this question, the article rests on three fundamental principles:

\begin{enumerate}[label=\textbf{\arabic*.}]
\item \textbf{Nature contains only one universal underlying medium.} What we call fields, forces, and particles are different phase, pressure, shear, rotational, or topological modes of this one substance.

\item \textbf{The stable unit of matter is an \emph{oscillon}}---a locally confined, self-sustaining resonant form of the underlying medium. Its intense internal dynamics leave a persistent pressure-deficit imprint in the surrounding medium.

\item \textbf{Gravity is the refractive response of the underlying medium.} What we perceive as attraction is, in fact, the wave-like motion of an object through an effective geometry modified by the pressure deficit.
\end{enumerate}

The Michelson--Morley experiment of 1887 detected no ether wind of the kind expected from the Earth's motion and placed stringent constraints on a stationary, matter-independent luminiferous ether \citep{michelson1887}. The underlying medium of RefG belongs to a different ontological picture: space, time, measuring instruments, and the observer are all modes of this single medium. RefG associates uniform motion with the symmetric wave-like, phase-coherent transmission of form, which leaves no dissipative wake on the corresponding Lorentz-covariant branch. This vision revives and extends Lord Kelvin's deep intuition: matter is itself a form of the medium, and the underlying medium is the substance from which the Universe is made.

The present text develops the ontological picture of RefG, giving intuitive physical meaning to the computational results collected in the public Zenodo archive. It connects the underlying medium with the origin of measurability, the pressure deficit associated with mass with cosmological horizons, and the origin of particles with the topology of quantum numbers. Medium-based models of unification have appeared before, including Winterberg's Planck aether \citep{winterberg2002}. Within this line of thought, RefG places particular emphasis on the operational origin of measurability and the mechanism of population-wide resonant locking. The following is a non-exhaustive list of the principal contributions of this work:

\begin{enumerate}
\item \textbf{Population-locked time.} The local rate of time, $\Omega(x)$, is generated by mutual locking within the particle population itself, $\nu_i=n_i\Omega(x)$; it does not rely on an externally prescribed metronome. Time dilation is a collective reconfiguration of this local network. This population-based definition of the temporal rate is a distinctive element of RefG.

\item \textbf{Two-channel accounting.} Gravitational mass is, by definition, a channel through which the medium is read externally. The internal state of the oscillon is retained as a physically separate quantity, and RefG separates the two axiomatically.

\item \textbf{Separation of extensive and intensive properties.} On the selected first-order biconformal branch, operational quantities scale according to a $p/p^2$ law. RefG relates these powers to the ontological distinction between nodes and their connections.

\item \textbf{The Koide $C_3$ phase structure.} In the three-phase map of the charged leptons, $\theta=h/9$ is the reduced coordinate of an order-nine return cycle. The $h=2$ branch gives $\theta=2/9$, while equalization of the internal channels yields $A_K=\sqrt{2}$.

\item \textbf{Effective closure of the fine-structure constant $\alpha$.} The constant $\alpha$ is interpreted as the ratio between internal phase--topological locking and external electromagnetic readout; an effective boundary calculation gives this relation a numerical form.

\item \textbf{Process-time quarantine.} A clear boundary is drawn between two temporal languages: coordinate time and process time. Process time remains an interpretive layer and does not enter the primary calculations of cosmological observables.

\item \textbf{Two readouts of a single operator.} The unified spectral hypothesis of RefG associates short modes with a particle map and long modes with a cosmic-web map.
\end{enumerate}

In this picture, dark energy is associated with the residue of the underlying medium's historical relaxation, while the three-dimensionality of space is associated with a candidate mechanism of knot closure and topological stability.

Modern tools for information processing, symbolic manipulation, and computation support systematic testing of the intuitive ideas developed in this research.

The text invites the reader to follow this chain of reasoning and see how a single physical picture connects phenomena usually treated separately. The scientific value of the framework will be determined by its further formalization and observational testing.

\section{Ontological Foundation}

\subsection{A Single Underlying Medium: Substance, Emptiness, and the Dynamics of Existence}

Modern fundamental physics describes the Universe in terms of many independent fields. In the Standard Model, each fundamental field has its own degrees of freedom and rules of interaction. Refractive Gravity (RefG) replaces this picture with a single foundational ontological principle: nature contains only one universal underlying medium. What we perceive as distinct fields and forces are different topological, phase, or rotational modes of this one substance. This monistic approach has deep roots in the history of physics. As early as 1870, William Clifford proposed that matter consists of small, moving ``hills'' of spatial curvature \citep{clifford1876}. Yet there is an essential difference between Clifford's picture and the RefG framework: the substrate in Clifford's picture is purely geometric, whereas the underlying medium of RefG supports pressure, shear, and resonant channels in addition to its geometric response.

When the Universe is conceived as a single medium, even a minimal picture of its dynamics requires two asymptotic limits. At one pole lies the absolute maximum of substance: a state in which pressure, mass--energy potential, and operational density approach their limiting values. At the other lies emptiness, or the zero limit, where physical interaction vanishes. This ``emptiness'' is a purely logical limit at which physical measurability falls to zero. Objective reality is the continuous, bounded oscillation between these two poles. The minimal active units of the underlying medium---here termed ``existence nodes''---emerge within this continuous dynamical process.

The ontological unity of the underlying medium comes with structural richness. It has several independent channels: it can be compressed or twisted, sustain vortex flow, and transmit wave excitations. Water offers a useful intuition. A single fluid can simultaneously carry hydrodynamic pressure, bulk flow, and surface tension, even though these properties cannot be reduced directly to one another. Likewise, the universal underlying medium combines phase, pressure, shear, and topological responses. This multiplicity of channels provides the physical resources for the formation of electric charge, spin, and a complex particle spectrum.

The inactive state of the underlying medium, denoted formally by $\varphi=0$, does not signify its ontological absence. The medium is ubiquitous; what varies is its local degree of excitation, the density of distinguishable nodes, and the strength of the coupling between them. In this language, a drop in pressure, or emptiness, corresponds to an increase in operationally unmeasurable intervals.

Here ``pressure'' denotes the effective pressure--energy variable of RefG. The condition $\varphi=0$ denotes the background state, while $\varphi<0$ denotes the deficit branch. Classical hydrodynamic pressure provides an intuitive analog of this variable. In the external readout, a change in $\varphi$ appears as a change in the local rate of time, the coordinate propagation of signals, the operational length scale, and the effective mass.

\subsection{The Origin of Measurability and Operational Space}

As long as the underlying medium remains in a state of absolute, inactive uniformity, measuring it is physically impossible in principle. In such a state, the concepts of ``distance'' and ``time'' do not exist, because there is no second physically distinguishable point for comparison. Physics proper begins with the medium's self-differentiation, or self-distinction. Once a local asymmetry---a physically distinguishable node, a phase displacement, or a pressure deficit---emerges within it, the medium is no longer a trivial, indistinguishable background.

The interval that arises between any two distinguishable nodes is the physical seed of space; the delay in the connection between them is the seed of time; and the energetic tension associated with differentiation is the ontological seed of mass. External gravitational mass is defined at a later stage through the readout of the deficit profile generated by this tension.

Physical meaning requires comparison. Three elements appear at the ontological origin: a \textbf{node}, a locally distinguishable state of the underlying medium; an \textbf{imprint}, the change that the node leaves in the surrounding medium; and a \textbf{relation}, the possibility of comparing one state with another. From this triad emerges the operational network that, in macroscopic language, we call space, time, and matter.

This picture is central to understanding the nature of the observer. We, like our measuring instruments, are complex composite structures built from the same nodal network. This approach strongly echoes the classical tradition of operationalism, especially Bridgman's principle that a physical concept is defined by the operations used to measure it \citep{bridgman1927}. Refractive Gravity goes further: it seeks the fundamental ontological condition that makes those measuring operations possible in the first place. In this account, the measurable Universe is not matter placed within some abstract ``space''; it is the continuous process by which the underlying medium distinguishes itself. Our operational space arises from the imprint of this distinguishability. Beyond that imprint, the universal medium may continue to exist, but it appears to us as a physical object only where it leaves a measurable, operational consequence.

The ontological separation between the nodes of the underlying medium is not measured directly. The observer is formed from modes of this same medium and has no measuring device outside it. What can be measured is the operational imprint of the interval between nodes. Operational space appears when the relation between at least two nodes becomes physically distinguishable. Below this threshold, the difference remains a local internal state of the underlying medium; once the threshold is crossed, it is read through changes in the local rate of time, operational length scale, mass, and signal propagation.

Differentiation does not acquire physical meaning merely by taking form; it does so when it crosses a minimum operational threshold and the relation becomes a measurable imprint. This threshold of distinguishability must lie somewhere, and its most natural physical candidate is a scale of Planck order. If the operational threshold does coincide with this scale, then a relation between nodes becomes physical reality precisely when that threshold is crossed. Any variation below it remains an internal microscopic vibration of the underlying medium and is not perceived by an observer as an independent spatial separation.

It follows from this principle that the ontological origin of the Universe is not a multiplicity of points. If an operational interval between nodes has not yet formed, then two physically distinguishable nodes do not yet exist, and neither does a spatial relation between them. Such a state cannot even be called a geometric ``single point,'' because the very concept of a point presupposes that space has already formed. A more precise term is an absolutely undifferentiated state of the underlying medium. Multiplicity appears in nature only when a difference arising in the medium leaves a stable physical imprint and the relation between nodes becomes operationally distinguishable. From this differentiation emerge the familiar structures of space, time, and the local rhythm of motion.

It is precisely at this fundamental boundary that the primary physical content of gravity comes into view. When the local connection between nodes of the underlying medium weakens, this ontological state appears in the external observer's readout as a reduction in pressure, a slowing of clocks, and a change in the conditions of signal propagation. A macroscopic observer perceives this coordinated ensemble of physical properties as a ``gravitational field.'' The following chapters develop this conceptual sketch into a quantitative refractive chain.

\subsection{The Local Collective Rhythm and Resonant Locking}

The shared rhythm of local time is not a preset metronome frequency imposed on the underlying medium from outside. It arises naturally when stable oscillons---the wave-like forms of matter---join in a common resonant network and their intrinsic frequencies settle into fixed, stable ratios to one another. In this collective state, the rate of time is not determined by an external background; it is created directly by the internal coordination and mutual locking of the local network itself.

At a deeper ontological level, the underlying medium is neither a tone nor a metronome. It is the instrument: a three-dimensional vibrating foundation that sets the boundaries, provides the resources, and determines the admissible vibrational spectrum. A bell offers a vivid analogy. Its sound is not produced by isolated notes; the modes of a single body resonate together. Likewise, the frequencies of elementary particles are established through mutual balance and coordination within the resources of the medium. The universal medium determines which modes can exist, while the particle modes balance one another within that admissible space and form a stable, locked configuration.

Denoting this shared local rhythm by $\Omega(x)$, the frequencies of stable oscillons may be written as

\begin{equation}
\nu_i(x)=n_i\,\Omega(x).
\label{eq:population-locking}
\end{equation}

Here $n_i$ is the fixed relative factor of the mode in question. In general, it is not an integer, so the particle frequencies do not reduce to simple harmonic overtones. They more closely resemble the mildly anharmonic spectrum of a complex three-dimensional volume resonator. We shall return to the three-phase structure of the charged leptons in later sections.

The relation $\nu_i=n_i\Omega$ leaves each mode with its own frequency while tying them all to one local collective scale. The ratios $n_i/n_j$ are therefore preserved under resonant locking. A rhythm incompatible with the network loses phase coherence and cannot persist as a stable particle mode.

In a low-pressure region, the collective scale $\Omega(x)$ of the entire local resonant network changes. In a dense medium, coupling is strong and $\Omega$ is high; in a rarefied medium, coupling weakens and $\Omega$ falls. Slower transfer between nodes, a lower energetic state of the medium, and a decrease in $\Omega$ are three readouts of the same state; each enters the quantitative chain only once.

When $\Omega(x)$ changes, every local clock and particle process scales with it. A local observer reads the rate of their own clock as normal. Under this common scaling, dimensionless ratios---including $\alpha$ and particle-mass ratios---remain invariant, because they are structural ratios of the form $n_i/n_j$.

The classical antecedent of this relational picture is Mach's principle \citep{mach1883}. RefG gives it a local form: the common rate of time is set directly by the local resonant network. The relational description is preserved, and physical action remains local.

In the effective $\alpha$ map of RefG, the fine-structure constant is interpreted as the ratio between internal resonant locking and external electromagnetic readout.

Now consider a wave imprint propagating from one pressure regime of the underlying medium into another. Its observed frequency is determined by the effective metric. In the ontological language of RefG, the comparison between the local rates of the source and receiver is read as a phase translation; cosmological redshift is defined by the metric relation $1+z=a_0/a_{\mathrm{em}}$.

Thus, within RefG, the rate of time is not set by an abstract universal clock outside local physics. It is the ontological state of the local resonant network itself. When a change in the pressure of the underlying medium transposes the network by a common factor, all compatible frequencies change together, while their dimensionless ratios remain invariant. This collective resonant locking is the essential physical content of the shared local rhythm and can be tested experimentally.

\subsection{The Oscillon, the Pressure Deficit, and the Origin of Mass}

What we conventionally call a ``particle''---whether an electron or a proton---is treated in the present framework as an oscillon. An oscillon is a locally confined, stable wave structure in the universal underlying medium. It is not an isolated rigid body independent of the medium; on the contrary, it is a highly organized, self-sustaining state of that medium. The concept has an interesting historical antecedent in Lord Kelvin's model of vortex atoms \citep{thomson1867}. Although that particular nineteenth-century program has passed into history, its conceptual core---the particle as a self-sustaining dynamical form of a medium---is the deep intuition that the oscillon picture revives and carries forward in modern resonant and topological language.

When an oscillon vibrates or rotates intensely, or stores energy within its form, the static pressure of the surrounding medium falls. Bernoulli's picture makes this direction of change vivid \citep{bernoulli1738}: an increase in dynamical intensity is accompanied by a decrease in static pressure. On the selected weak, static branch of RefG, this analogy corresponds to a pressure--energy balance identity.

By existing as an active form, an oscillon generates a pressure-deficit profile in the underlying medium. In the external gravitational readout, this profile is the source of mass: the more intense the oscillon's internal dynamics, the deeper the deficit and the stronger the refractive gradient.

This picture has a well-known historical parallel in John Archibald Wheeler's program of ``mass without mass,'' in which a geon is a massive configuration held together by the energy of its own field \citep{wheeler1955}. Wheeler's geons gradually radiate their energy away; RefG associates the stability of an oscillon with resonant locking and topological protection.

The term \emph{effective refractive mass} denotes the external readout of an individual oscillon in a given background. Its internal state---its topological structure and internal energy---is invariant in the local description, while the external gravitational readout changes with the background. The asymptotic mass of the complete system belongs to a separate integral readout channel and does not follow automatically from the $m_{\mathrm{eff}}$ of an individual oscillon. With this distinction, the origin of mass is read as the formation of the medium's refractive response.

At this fundamental level, gravity is the refractive deflection of one oscillon's wave form within the pressure deficit already created in the medium by another. Here, a gravitational well is not a mechanical ``stretching'' or deformation of the medium, but a reduction in its local pressure and energetic state. This deficit is the physical source that changes the medium's effective geometry, and the particle's spatial motion then unfolds within that altered geometry.

An oscillon generates a pressure deficit, and the deficit in turn changes its effective mass, operational size, and coupling to the medium. RefG reads this two-way feedback as a mechanism of dynamical equilibrium: the material form changes the medium, while the altered medium determines the conditions of the form's external physical manifestation.

\section{Gravity as Refraction}

\subsection{Three Concurrent Operational Changes}

As noted in the preceding chapter, the interval between nodes of the underlying medium is not measured directly; only its physical imprint is operationally accessible. Around an oscillon, this imprint appears as a local reduction in the pressure state of the medium: physically distinguishable nodes are more sparsely distributed, and the informational connection between them is weaker. For an external macroscopic observer, this single ontological state manifests itself through three simultaneous, coordinated changes.

\paragraph{First: a reduction in the coordinate speed of light.}
Within a pressure-deficit region, nodes are more sparsely distributed and their mutual interaction is weaker. The propagation of light and other wave signals therefore slows in the coordinate description. The locally measured speed of light remains invariant; what changes is the propagation rate expressed in the coordinates of a distant observer.

\paragraph{Second: a reduction in the operational length scale of the oscillon.}
The fundamental topological form of an oscillon does not undergo mechanical compression. Rather, in a rarefied underlying medium, that form contains fewer interacting nodes. In the external description, this is read directly as a smaller operational size of the oscillon and, correspondingly, as a modified effective-mass readout.

\paragraph{Third: a reduction in the rate of the local resonant clock.}
A clock is a real physical system built from oscillons, and its rate depends on local nodal connections. In a rarefied medium, those connections weaken, and the clock runs more slowly in the coordinate time of a distant observer. The local clock itself continues to register its own rhythm as normal.

A double-walled vacuum flask offers a vivid analogy: fewer carrier molecules mean fewer collisions and slower heat transfer. In the RefG map, this corresponds to the relation between the strength of nodal coupling and the coordinate rate of a clock.

These three changes are operational readouts of the same reduced pressure state. They do not, however, depend on the change in the medium in the same quantitative way.

\begin{figure}[htbp]
\centering
\includegraphics[width=0.92\linewidth]{figure_3.pdf}
\caption{Operational readout of an oscillon in a pressure-deficit region. \textbf{(A)} In a high-pressure underlying medium, the oscillon occupies a dense nodal configuration; inter-node coupling is strong, and local processes proceed at a high rate. \textbf{(B)} In a pressure-deficit region, the nodes of the underlying medium are more sparsely distributed; the operational separation between them increases, the coupling rate falls, and local time slows. \textbf{(C)} In the external readout, an oscillon within the pressure deficit is read with a slower temporal rate, a smaller operational size, and a modified effective mass; the enlarged internal interval does not enter directly as a measurable quantity.}
\label{fig:oscillon-pressure-deficit}
\end{figure}

The separation of nodes does not generate a ``second space.'' The interval between nodes is not an independent physical object; within this framework, ordinary ``emptiness'' is read through the pressure variable. The operational manifestation of a rarefied underlying medium is a local reduction in its pressure state.

\subsection{The Quadratic Relation and the Factor of Two}

The three operational changes do not scale with the same power. In this map, the operational size and effective mass of an individual oscillon belong to the node-count channel, whereas the coordinate speed of light belongs to the channel that governs the strength of inter-node coupling. The first readout is extensive; the second is intensive, meaning that it does not depend on the size of the system.

Writing the pressure factor as $p$, the scaling takes the form

\begin{equation}
\frac{L_{\mathrm{oper}}}{L_0}
=\frac{m_{\mathrm{eff}}}{m_0}
=\frac{\Omega}{\Omega_0}
=p,
\qquad
\frac{c_{\mathrm{coord}}}{c_0}=p^2,
\qquad
p=e^{\varphi/2}.
\label{eq:operational-scaling}
\end{equation}

The same powers also appear in polarizable-vacuum models \citep{dicke1957,puthoff2002}. RefG relates them to the ontological distinction between nodes and connections.

In formal terms, the external macroscopic readout of an individual oscillon is related to the pressure-deficit profile of the underlying medium by a specific exponential half-factor:

\begin{equation}
m_{\mathrm{eff}}=m_0 e^{\varphi/2}.
\label{eq:effective-mass}
\end{equation}

On the deficit branch, $\varphi<0$, and hence $0<p<1$. This expression describes the external readout of an individual oscillon placed in a given background: as $\varphi$ decreases, its operational size and $m_{\mathrm{eff}}$ decrease. The total gravitational charge of the system that generates the deficit is defined separately through an integral readout. This relation follows from the metric filters and the selected first-order biconformal branch of the theory.

Four quantities must be distinguished:

\begin{enumerate}
\item the oscillon's internal state and its own energy;
\item the effective mass $m_{\mathrm{eff}}$ of an individual oscillon in a specified background;
\item the profile of the deficit source generated in the underlying medium; and
\item the ADM/Komar-type asymptotic gravitational charge of the complete system \citep{arnowitt1962,komar1959}.
\end{enumerate}

On the selected static branch, the projection of the deficit profile is read as an effective four-component source for the metric. A similar distinction exists in standard gravity: the local energy of a test body and the integrated mass of the entire system are not the same quantity.

A vivid analogy helps us picture this complex, self-consistent dynamics. Imagine a colored, ice-like crystal immersed in clear water. The crystal is a locally bound form of the same water, just as an oscillon is a local resonance of the common underlying medium. As the crystal dissolves, its color passes into the water: the surrounding medium gradually becomes colored, while the crystal itself shrinks. Energy does not disappear. Like the color dissolving into the water, it passes from a locally bound state into a broader state of the medium.

As more of the medium becomes ``colored'' by this imprint, its local pressure state falls, and the externally read mass, operational size, time, and distance change. This is a relaxation process: the more of its imprint the oscillon transfers to the surrounding medium, the less capacity it retains to transfer any further imprint. The process therefore slows, and the deepening of the deficit is limited by its own accumulated trace.

Against an external measurement standard, this evolution approaches the minimum resolvable scale---possibly of Planck order---where the object appears to vanish from the external readout. Local measuring instruments participate in the same relaxation regime, so this limit arises differently in the internal readout. The mass deficit and the energy transferred to the underlying medium are two independent readings of a single process.

This scaling has the biconformal metric representation

\begin{equation}
ds^2=e^\varphi c^2dt^2-e^{-\varphi}d\boldsymbol{\sigma}^2,
\qquad
\varphi=-\frac{2U}{c^2},
\label{eq:bi-conformal-metric}
\end{equation}

where $U>0$ is the potential of the selected exponential branch and, at leading order, agrees with the magnitude of the Newtonian potential. In the weak-field regime,

\begin{equation}
g_{tt}=1-\frac{2U}{c^2}+\frac{2U^2}{c^4}+\cdots,
\qquad
g_{ij}=-\left(1+\frac{2U}{c^2}+\cdots\right)\delta_{ij}.
\label{eq:weak-field-expansion}
\end{equation}

On the selected first post-Newtonian (1PN) weak-field branch, these expressions give the parametrized post-Newtonian (PPN) coefficients $\beta=1$ and $\gamma=1$.

The duration of one complete clock cycle is given by the ratio of operational size to coordinate speed. If the size is halved while the coordinate speed falls by a factor of four, the duration of the cycle doubles: the clock runs at half its original rate.

On the same branch, the first post-Newtonian results agree with the standard values for light deflection, Shapiro delay, gravitational redshift, and perihelion precession \citep{bertotti2003,will2014}. Historically, the additional factor of two supplied by the spatial component completes Einstein's 1911 half-value to the full general-relativistic deflection \citep{einstein1916,shapiro1964}. The higher-order Solar System branch is selected by a separate field equation and does not follow automatically from the exponential representation above.

\subsection{Refractive Attraction, Inertia, and the Absence of Friction}

An oscillon is a wave structure embedded in gradients of pressure and coordinate propagation speed across the underlying medium. Its phase propagates faster on the dense side and more slowly on the rarefied side, causing the wavefront to refract toward the region of lower pressure. In the weak, static, large-scale limit, the deficit generated by an object appears as an effective geometry, while wave propagation through that geometry appears as gravitational motion.

RefG describes a falling apple as follows: the wave forms of its atoms bend toward the region of lower pressure and slower coordinate propagation created by the Earth. Familiar gravitational attraction is the refractive response of the underlying medium.

In the formal description, the gravitational source is the pressure deficit projected into the underlying medium. It is the component of the pressure--energy balance that an external observer reads as attraction. Microscopically, it is the stress profile imprinted on the medium by the oscillon's internal dynamics.

A gravitational well is a three-dimensional region of reduced pressure in the underlying medium. The mass of a bound system is smaller than the sum of the masses of its freely separated constituents by the binding energy.

GW150914 provides an observational example: gravitational waves carried away energy equivalent to approximately three solar masses \citep{abbott2016}. RefG reads this event as a unified energy map of structural reconfiguration and tensor radiation. A complete dynamical calculation determines the quantitative relation between the source profile and the radiated transverse--traceless (TT) flux.

On the static compact branch, the projected deficit source appears as a violation of the effective null energy condition (NEC). In the RefG map, this stress belongs to the confinement channel of the underlying medium.

RefG associates inertia under acceleration with deformation of the deficit imprint. The pressure profile becomes asymmetric ahead of and behind the oscillon, and this asymmetry produces a reaction opposing the acceleration. Under deceleration, the configuration is reversed.

On the branch describing a localized, finite-energy, Lorentz-covariant oscillon, the leading inertial coefficient corresponds to the full Noether energy when the Laue boundary condition is satisfied \citep{noether1918,laue1911}. Under the same conditions, inertial and gravitational masses are read out together at leading order. The MICROSCOPE experiment tests this universality at the level of $10^{-15}$ \citep{touboul2022}.

On the uniformly moving Lorentz-covariant solution branch, the wave pattern is transmitted symmetrically and leaves no dissipative wake. RefG associates frictionless motion with this regime of phase transmission.

Motion is here a wave-like relay of form. Its limiting speed is set by the maximum rate of phase transfer between nodes of the underlying medium. The finite speed of light emerges from the same operational network.

\subsection{Local Unobservability and the Principle of Relativity}

A local observer in a gravitational field finds their own temporal rate and spatial scale entirely normal, because the observer's clock, ruler, and body are all composed of the same underlying medium. When the medium changes, the entire local system is reconfigured by a common factor. A change that is sufficiently slow relative to the internal periods is adiabatic: it preserves the mode without exciting unwanted transitions. Under this common scaling, dimensionless physical ratios, including $\alpha$, remain invariant.

This picture is subject to stringent experimental bounds. Gravitational redshift has been tested across several regimes, from the Pound--Rebka experiment to Gravity Probe A and the Galileo tests \citep{pound1960,vessot1980,delva2018,herrmann2018}. Quasar spectra and atomic clocks place independent constraints on the stability of dimensionless ratios. RefG interprets these observations as physical manifestations of local self-calibration.

If the underlying medium changes uniformly throughout a given region, rulers, clocks, and the conditions governing light propagation all change in concert, and the local readout self-calibrates. A gravitational wave, by contrast, is a time-dependent, quadrupolar differential deformation: it acts with different signs along different axes and therefore leaves a measurable imprint in a detector.

The FitzGerald--Lorentz length compensation provides a historical antecedent to this logic \citep{fitzgerald1889,lorentz1892}. RefG attributes the compensation to the fact that measuring instruments are themselves wave excitations of the underlying medium. In the local Minkowski limit, the tensor branch is luminal and produces a transverse--traceless detector response. A possible scalar breathing branch describes an additional volumetric deformation of the medium. These branches are constrained by the GWTC-3 tests of gravity \citep{ligovirgokagra2022}.

The clearest operational invariant in this ontological picture is the difference between the proper times accumulated by two clocks that follow different trajectories and later reunite. The numbers recorded on their displays are objective consequences of their histories through different gravitational regimes. Locally measured rates, sizes, and frequencies remain relational readouts of processes unfolding within the common medium.

\subsection{Historical Context and the Place of RefG}

The refractive account of gravity has clear antecedents in the history of physics. In 1911--1912, before the final formulation of general relativity, Albert Einstein related the coordinate speed of light to the Newtonian gravitational potential through an expression of the form $c=c_0(1+\Phi_N/c^2)$ \citep{einstein1911,einstein1912}. The optics of physical media reaches further back. In 1851, Hippolyte Fizeau showed that moving water partially ``drags'' the propagation of light \citep{fizeau1851}, while Walter Gordon later related light propagation in a dielectric medium to the geodesics of an effective metric \citep{gordon1923}. Together, these results established a formal bridge between the optics of media and gravitational geometry.

In 1957, Robert Dicke recast gravity in terms of a variable refractive index of the vacuum \citep{dicke1957}. A foundational result in modern analog gravity came from William Unruh, who showed that acoustic excitations in a moving fluid obey an effective metric analogous to a black-hole horizon \citep{unruh1981}. Volovik's condensed-medium program, together with the work of Barcel\'o, Liberati, and Visser, showed how an effective metric can emerge from the density and dynamics of a background medium \citep{barcelo2011}. This line of research also received a striking laboratory extension: in 2019, a thermal analog of Hawking radiation was observed in a Bose--Einstein condensate \citep{munozdenova2019}.

Alongside this physical-medium tradition stands a thermodynamic branch of emergent gravity. Andrei Sakharov interpreted gravity as an effect generated by quantum fluctuations of the vacuum \citep{sakharov1967}, while Ted Jacobson showed that the Einstein equation can be obtained as an equation of state from the thermodynamic conditions of local horizons \citep{jacobson1995}. This branch is kindred to RefG in its view that spacetime geometry is nonfundamental; in RefG, effective geometry is generated by a specific pressure--refractive medium.

Refractive Gravity adds three fundamental new elements to this rich historical line:

\begin{enumerate}
\item \textbf{An ontological foundation.} RefG introduces a single universal underlying medium that is not merely an arena for the propagation of light, but also gives rise to matter (oscillons), measuring instruments, and observers. This ontological step carries the approach from an auxiliary mathematical analog into an independent physical research framework.

\item \textbf{A local collective resonant rate.} The collective resonant rate of oscillons makes the local unobservability associated with the principle of relativity an intrinsic dynamical consequence of the system.

\item \textbf{A universal pressure-deficit channel.} RefG uses the pressure deficit as a common language for both source and readout. This language connects mass, temporal rate, light deflection, and the local tensor response; the complete source-to-flux bridge for gravitational waves belongs to a separate dynamical calculation.
\end{enumerate}

Within this historical map, the closest conceptual counterpart to RefG is Erik Verlinde's entropic gravity \citep{verlinde2011}. In Verlinde's model, Newton's law follows from the entropic balance of information on a holographic screen. RefG shares the idea that gravity is emergent but takes the pressure deficit of the underlying medium as its source. The two mechanisms predict different observational consequences and can be tested independently.

\section{Three Scales and Cosmological Horizons}

\subsection{The Local Scale: Compact Objects, Internal State, and External Readout}

The universal pressure chain developed in the preceding chapters manifests itself in nature on three entirely different astrophysical scales: locally, as the pressure deficit surrounding an oscillon; galactically, as a coherent flow of the universal underlying medium; and cosmologically, as a global historical change in the state of that medium. These are, in fact, three dynamical regimes of one and the same underlying medium. We begin at the local scale, where the chain takes its sharpest and most extreme form---in what are known as compact objects.

In the extreme-gravity regime, an object's \emph{internal state} becomes sharply distinct from its \emph{externally read mass}. The compact-object map in RefG separates the internal state of accumulated matter from the mass read at a distance. A distant observer reads that state through the deficit profile, the clock rate, and the operational size. The complete solution determines the internal geometry, the horizon, and the causal structure.

RefG relates the mass difference associated with a merger to changes in binding and in the deficit profile. Its quantitative relation to the radiated tensor flux is fixed by the complete dynamical solution. In Chapter II, GW150914 provided an observational example of this energy map.

The metric determines the difference between internal and external scales. On the compact branch of RefG, that difference can be large. The exact internal scale, the existence of a horizon, and the global causal structure follow from the complete solution.

Penrose's singularity theorem rests on several assumptions, including the relevant energy and causal conditions \citep{penrose1965}. On the compact branch of RefG, the effective source violates the null energy condition (NEC), and the theorem's conclusion therefore no longer follows automatically. A regular interior is determined by the complete field solution, its boundary matching, and its stability.

The idea of regular compact objects has also been developed in the modern literature. A well-known example is the ``gravastar'' of Mazur and Mottola, a condensate object with no singular core \citep{mazur2004}. Alternatives of this kind may be distinguished by their shadow size, oscillation spectrum, and echo signals \citep{cardoso2019}. RefG is structurally related to this family, while tying the interior regime to the pressure--energy balance of the universal underlying medium.

RefG connects the feedback between mass and deficit to a finite equilibrium state. The field equations determine whether this state exists and is stable. In this picture, a compact object is a heavily loaded regime of the underlying medium whose internal state and remotely read mass are registered through different channels.

\subsection{The Galactic Scale: Vortical Flow and MOND}

\paragraph{Observational law.}
Galaxy rotation curves do not fall off as the distribution of visible baryonic matter alone would suggest. At the same time, baryonic mass and rotation speed obey a tight relation---the baryonic Tully--Fisher relation. Historical landmarks are the rotation curve of Andromeda \citep{rubin1970} and the Tully--Fisher relation itself \citep{tully1977}. In the modern radial acceleration relation, 2,693 measurements from 153 galaxies lie on a single common curve \citep{mcgaugh2016}.

The Lambda cold dark matter model ($\Lambda$CDM) describes these phenomena with dark-matter halos. MOND instead singles out a characteristic acceleration scale, $a_0$, at which departures from the Newtonian expectation appear \citep{milgrom1983}. Tensor--vector--scalar theory (TeVeS) and its newer generalizations use additional vector and scalar fields to provide a relativistic description of this phenomenology \citep{bekenstein2004,skordis2021}. The test for RefG is clear: $a_0$ must be tied to the dynamics of the underlying medium.

\paragraph{Candidate vortical mechanism in RefG.}
The physical structure behind this map is a large-scale \textbf{vortical, or swirling, flow} of the underlying medium. An additional pressure deficit forms at the center of this ordered rotational regime, and baryonic matter gathers in the resulting well. The flattening of rotation curves and the Tully--Fisher relation are read as the combined refractive effect of a giant vortex in the underlying medium and the imprints of baryonic oscillons.

The closest modern relative of this picture is superfluid dark matter, in which a MOND-like force is associated with phonon excitations of a condensate phase \citep{berezhiani2015}. There, superfluidity is a phase of an additional particle component; in RefG, coherent flow is a regime of the single underlying medium itself.

\paragraph{Morphology.}
A galaxy's visible form may reflect the flow regime of the underlying medium. Ordered, laminar rotation favors disk and spiral structures; a more chaotic or isotropic flow favors dispersion-supported systems. These are the two endpoints of a continuous spectrum, since some elliptical galaxies also retain substantial rotation.

\paragraph{Clusters.}
Changing the acceleration law alone often fails to recover all the mass required in galaxy clusters; historically, the missing-mass problem first emerged in the Coma Cluster \citep{zwicky1933}. RefG reads a cluster as a deep node in a long-wavelength resonant network. The cumulative deficit of its galaxies and their extended oscillon imprints adds to this background deficit. In Chapter IV, particles and the cosmic web will appear as opposite ends of the same spectrum.

This picture suggests two candidate observational signatures.

\paragraph{First---quantization of circulation.}
If the flow of the underlying medium has a single-valued condensate phase, its expected unit of circulation is $h/m$, where $m$ is the mass of the effective flow mode. For an extremely small $m$, individual vortex imprints may appear in weak gravitational lensing or in disk kinematics. Similar signatures are sought in the wave-like halos of ultralight dark matter \citep{hu2000,hui2021}.

\paragraph{Second---the frozen imprint of the Bullet Cluster.}
In the Bullet Cluster, the centers inferred from weak gravitational lensing are offset from the X-ray plasma \citep{clowe2006}. RefG reads this separation as a measure of vortical memory retained in the underlying medium. A quantitative model must recover both the lensing map and the spatial distributions of the gas and galaxies.

The resulting test map is unified: it spans the galactic data for the SPARC radial acceleration relation (RAR) \citep{lelli2016,mcgaugh2016}, galaxy clusters, Bullet-Cluster-type mergers, the cosmic microwave background (CMB), and large-scale structure. On a branch without dark matter, the vortical source must also replace the cosmological contribution of cold dark matter.

\subsection{The Cosmological Scale: Relaxation of the Underlying Medium and Two Languages of Time}

At the cosmological scale, the external metric description remains unchanged: the Universe expands, distances between galaxies increase, and light undergoes cosmological redshift. Viewed from within, RefG reads this history as the relaxation of the underlying medium. The long-range imprints of oscillons accumulate into a common background and alter the medium's pressure state. In the language of process time, earlier epochs correspond to a larger internal budget of cycles for structure formation.

Photons follow null geodesics of the effective metric, and cosmological redshift obeys the standard scale-factor relation, $1+z=a_0/a_{\mathrm{em}}$. In the ontological language of RefG, this is the phase translation described in Chapter I: an imprint arriving from an earlier state of the underlying medium is read against the local resonant rate in each new epoch, while dimensionless ratios remain unchanged.

The ``age of the Universe''---approximately 13.8 billion years---is metric time calibrated by present-day clocks \citep{planck2020}. Before the first stable modes of matter formed, no operational clock had been defined. Process time provides an additional description of the internal history of the underlying medium.

The idea of two time scales in cosmology has historical antecedents. Milne described the same Universe in terms of a kinematic time $t$ and a dynamical time $\tau$ \citep{milne1935,milne1948}. Wetterich linked the languages of expansion and changing material standards within a specific field model \citep{wetterich2013}. In RefG, process time is an internal interpretive scale; in this language, the rapid maturation of early galaxies appears as a possible sign of an additional budget for structure formation \citep{labbe2023}.

Cosmological expansion is expressed in RefG at three descriptive levels:

\begin{enumerate}
\item \textbf{In the external metric language,} distances in space increase.
\item \textbf{In the material language,} clocks and atomic standards from different epochs are situated in different local states of the underlying medium.
\item \textbf{In the language of the underlying medium,} the mean density of nodes and the connections between them change globally over the course of cosmic history.
\end{enumerate}

\begin{quote}
These are three levels on which the same history can be described. The external metric is the primary language of observational calculations.
\end{quote}

The limit of the medium's relaxation is defined thermodynamically: diminishing gradients mark an approach to maximum entropy. RefG reads this process as an \textbf{asymptotic} evolution, with equilibrium remaining a limiting state.

RefG associates the macroscopic arrow of time with the direction in which the underlying medium relaxes. The dissolving-crystal picture from Chapter II makes this intuition tangible: color dispersed through water does not spontaneously gather back into a solid crystal. RefG reads this statistical irreversibility as the macroscopic direction of time.

In this interpretive map, the thermodynamic, cosmological, and radiative arrows share a statistical source in the relaxation of the underlying medium. Macroscopic irreversibility reflects the one-way approach to equilibrium \citep{planck1899}, while the arrow of time is tied to the medium's initial nonequilibrium state.

The late Universe is expanding at an accelerated rate. Standard cosmology attributes this phenomenon to $\Lambda$ or dark energy. RefG reads it as an energetic residue of the underlying medium's relaxation. The stress of this residue gives rise to negative pressure, while its slow variation yields an almost constant cosmological term. The numerical value of $\Lambda$ must emerge from the relaxation equations.

Within the ontology of a single medium, the vacuum-energy problem becomes a question of the equilibrium stress--energy tensor of the underlying medium and its gravitational readout. RefG reads the small cosmological residue as a residual departure from equilibrium. Related mechanisms have been studied in condensed-matter physics and in $q$-theory \citep{volovik2003,klinkhamer2008}. The same underlying medium gives rise to both gravitational refraction and late-time relaxation. The quantitative curve must connect $\rho(a)$, $w(a)$, and $\Lambda$, and must be tested against baryon acoustic oscillations (BAO), Type Ia supernovae, and CMB data \citep{abdulkarim2025,brout2022,planck2020}.

\paragraph{Process-time quarantine.}
Process time is an internal interpretive scale. The expansion history, the CMB, Big Bang nucleosynthesis (BBN), photon redshift, atomic clocks, and stellar nuclear processes are calculated in metric time. The spectral evolution of supernovae follows the observed $(1+z)$ time dilation \citep{blondin2008}. The process-time scale does not enter the formulae for these primary observational calculations.

\subsection{Quantum Randomness and Two-Level Time}

The twofold nature of time also opens an interesting interpretive possibility in quantum physics. The place and moment of a measurable event are registered in oscillon time---that is, clock time. RefG associates the selection of an outcome with the internal level of the underlying medium. Clock time does not measure this level directly, because the clock itself is an oscillon structure produced by the same medium.

Violations of Bell inequalities rule out local hidden-variable models that explain outcomes through locally preassigned parameters \citep{bell1964,hensen2015}. At the same time, the no-signaling condition preserves relativistic causality: distant correlations cannot be used to transmit controllable information faster than light.

RefG reads this coexistence as \textbf{atemporal internal coordination}. In the language of the underlying medium, it is not a signal traveling through space; in the language of clocks, it appears only as a correlation between distant outcomes.

RefG reads Wheeler's delayed-choice experiment in the same language \citep{wheeler1978}. In a single-photon realization, the measurement arrangement chosen at a late stage determines which complementary outcome appears: interference in a closed interferometer, or path information in an open one \citep{jacques2007}. RefG's two-level time describes these facts without a retrocausal signal.

This picture has influential relatives in modern theoretical physics: the timeless fundamental level of the Wheeler--DeWitt equation; the Page--Wootters mechanism, in which time emerges from internal correlations among subsystems; and Bohm's pilot-wave picture, with its nonlocal hidden coordination \citep{dewitt1967,pagewootters1983,bohm1952}. RefG joins this family with an ontology of its own, grounded in the underlying medium.

A quantitative completion of this quantum picture must satisfy two stringent conditions:

\begin{enumerate}
\item it must recover the \textbf{Born probability distribution} as the macroscopic statistical readout of internal coordination \citep{born1926}; and
\item it must preserve the \textbf{no-signaling condition}: internal coordination cannot become a transmissible signal that violates relativistic causality.
\end{enumerate}

\section{Oscillons and the Origin of Matter}

\subsection{Many Patterns in One Ocean}

For a vivid picture, imagine an unbounded, continuous ocean with waves rising within it. Most of them form, propagate, and soon disappear. A few, however, settle into stable, confined patterns: one may take the form of a ring, with the flow circulating continuously as a vortex; another may be a stable spherical pulsation; a third, a spiral wave.

\begin{figure}[htbp]
\centering
\includegraphics[width=0.95\linewidth]{ai-render-16799135.jpg}
\caption{Resonant modes of the underlying medium. In the intuitive map of RefG, stable wave patterns in a single medium appear as different particle modes: a spherical pulsation, a ring-shaped vortex, and a spiral wave. Here a particle is not a separate substance, but a locally confined resonant pattern of the underlying medium.}
\label{fig:resonant-modes}
\end{figure}

The ocean picture conveys the central proposition of this chapter: the electron, proton, photon, and quark are different resonant patterns of one universal underlying medium. The medium itself is what oscillates; the particle is a stable form of that oscillation. A historical precursor is the Skyrme model, in which nucleons are topological solitons of a nonlinear meson field \citep{skyrme1961}. RefG takes the intuition one step further: every such form is a mode of the same underlying medium.

The classical image of a small, solid sphere gives way here to a membrane-like picture: the vibrational amplitude of the underlying medium extends through space, while a localized event in a detector is the resonant contact between that vibration and the detector's internal oscillon modes. In this way, RefG gives the spatial amplitude and phase of the quantum state an ontological foundation in the underlying medium.

Single-electron interference brings two facts together: the state's amplitude depends on the geometry of both slits, yet each detection occurs at a single point on the screen. Tonomura's group directly recorded the gradual build-up of this pattern \citep{tonomura1989}. Likewise, a stationary electronic state in an atom does not radiate as a classical orbit would. RefG reads these results as the conjunction of an extended wave in the underlying medium and a localized detection event.

In high-energy scattering, the electron appears pointlike in its electromagnetic form factor. A Compton-scale oscillon must therefore reproduce this behavior across the accessible momentum range. The form factor derived from the oscillon's charge profile is a direct test of that requirement.

\subsection{Localized Confinement: Massive and Massless Modes}

In RefG's particle map, the principal distinguishing feature between massive and massless modes is \textbf{localized confinement}.

In this map, the energy of a massive oscillon is confined within a localized form. When vortical circulation or stable pulsation closes a wave pattern within a finite region of space, that pattern creates a pressure-deficit profile. Long-lived localized pulsations are known in nonlinear field theory as pulsons and oscillons \citep{bogolyubsky1976,gleiser1994,copeland1995}. RefG ties their stability to resonant locking across the population. In the external readout, the confined energy appears as inertia and gravitational mass.

A massless carrier mode has no rest frame of its own and no stable localized core; it propagates through space at the limiting speed. RefG reads the photon as such a propagating phase mode.

Helicity and confined local rotation are not the same. The photon's helicity $\pm1$ describes the helical structure of phase along the direction of propagation; it does not create vortical circulation confined to one place. A massive mode has a stable localized profile, whereas a massless carrier has a propagating dynamical phase mode.

A photon's energy and momentum enter the gravitational balance. Free radiation leaves no persistent, localized deficit profile of the kind associated with a massive oscillon, but its energy contributes in full to the source of the complete system.

The same rule extends to composite particles. Strong-interaction energy confined within a hadron is not read as a separate external imprint, but it contributes to the hadron's total energy and gravitational mass.

\subsection{Self-Forming Resonators and the Two-Level Filter}

A wave equation permits an effectively infinite variety of forms, yet nature presents us with a sparse, well-defined spectrum of stable particles. Why does nature select this particular set from infinitely many possibilities? RefG's answer rests on a \textbf{two-level filter}.

The first level is a kinematic and population filter. A handpan provides the intuition. A guitar predetermines the resonant length of its string; a three-dimensional oscillon, by contrast, must ``forge'' its own resonant cavity within the medium. In a stable handpan note, the fundamental, octave, and fifth come together in one local structure. In the RefG map, an oscillon's deficit imprint changes the local geometry, and the altered geometry in turn selects the viable modes. A mode must be able to propagate through the underlying medium and coexist with the existing particle spectrum. Stability emerges from compatibility across the entire local population.

The second level is a filter for energetic stability. A form that passes the first filter but has an open kinematic decay channel decays into lower-energy modes. Hadronic resonances of the $\Delta$, $\rho$, and $\omega$ type are familiar examples. Electron decay has not been observed, and charge conservation strongly protects the electron as the lightest charged particle. Proton decay has not been observed either, although its absolute stability requires protection by a separate symmetry.

The mathematical language of this filter is that of \textbf{attractors} in nonlinear dynamics---stable states toward which a system naturally evolves. The energy window selects which compatible modes survive as long-lived particle states.

A stable particle is a self-consistent equilibrium: the wave pattern creates a deficit in the medium, and through feedback that deficit selects and reinforces the pattern itself. In nonlinear fields, Q-balls show how a conserved internal charge can protect a localized configuration \citep{coleman1985}. In RefG, phase locking across the population also plays a protective role. The gravitational feedback between mass and deficit is read as the microscopic continuation of the same equilibrium.

The underlying medium is the common vibrational channel, while phase locking occurs directly among the oscillons. They synchronize with the existing spectrum through the fixed resonant ratios described in Chapter I.

This mutual locking requires neither a master clock nor an external conductor. Huygens's pendulums suspended from a common beam, phase-locked laser modes, and the Kuramoto model are familiar examples of spontaneous synchronization \citep{huygens1665,kuramoto1975}. A historical relative is the bootstrap program in particle physics, where the spectrum is understood as an ensemble of mutually compatible modes \citep{chew1961}.

A sufficiently slow change in the gravitational field preserves these locks. The change is adiabatic when it proceeds slowly enough relative to the internal periods to avoid exciting unwanted transitions. In a rapidly changing, strongly nonadiabatic background, the mismatch with the surrounding medium may exceed the binding strength and break the lock.

When a wave imprint enters a different pressure environment, it is read anew under the local operational conditions: its wavelength, transfer rate, and frequency relative to local clocks change, while the phase ratio is preserved. This is the transposition principle described in Chapter I.

\subsection{The Global Resonator: Unity of the Microscopic and Macroscopic Spectra}

The two-level filter tells us \emph{how} the particle spectrum is selected. We now ask: where is it realized? Ernst Chladni's plate provides the intuition \citep{chladni1787}. Sand scattered across a vibrating metal plate gathers along the nodal lines of a standing wave, where the motion is minimal. In the RefG map, this picture corresponds to global modes and their nodes. Every resonator has a spectrum of its own: geometry and boundaries determine which modes are amplified and which fade away.

In RefG's global-resonator hypothesis, the Universe is a \textbf{global volumetric resonator}---a three-dimensional underlying medium with its own allowed spectrum. The strong excitation at the cosmological beginning activates that spectrum as a whole. Modes reinforce or cancel one another; their deficit imprints lock into resonance or drift out of phase. A form incompatible with the geometry and background pressure cannot remain a stable particle, and its energy passes into the allowed channels.

In this hypothesis, today's stable particles are the local, short-wavelength modes of the universal spectrum. Where phase ordering did not have time to complete in the early Universe, frozen lines of topological defects may have remained in the underlying medium. We return to this cosmological imprint in Chapter V.

In the same map, dispersion in the underlying medium links frequency to spatial wavelength. Short modes correspond to localized particles; long modes, to the large-scale distribution of matter---the cosmic web.

Galaxies here resemble Chladni's sand: their distribution reveals the stable nodes, filaments, and vast voids of the underlying medium. Galaxy clusters are the deepest nodes in this network, where long-wavelength resonance lowers the background pressure. This is the same picture of clusters developed in Chapter III.

The conclusion is both simple and beautiful: the microscopic particle and the macroscopic cosmic web are two readouts of one universal spectrum---a form bound locally on the inside, and a pattern spread globally on the outside. In standard cosmology, the filaments, sheets, and voids of the cosmic web are described by Zel'dovich's anisotropic collapse \citep{zeldovich1970}. RefG expresses this morphology in a resonant ontological language.

RefG's candidate map of a unified spectral operator associates short modes with the particle readout and long modes with the background nodal deficit. The observed BAO scale remains a standard ruler \citep{eisenstein2005}; RefG reads it as the cosmological imprint of the underlying medium's long-wavelength spectrum.

A finite, boundaryless resonator of the underlying medium is a cosmological hypothesis of RefG. Wetterich's model shows that a change of variables can recast expansion as a relation between geometry and material standards \citep{wetterich2013}. In RefG, the compatibility of these languages is tested through the agreement of all dimensionless observables.

In the finite-resonator hypothesis, the size of the underlying medium is fixed, while cosmological history is translated through changes in internal standards and operational readout. The CMB, BBN, and $T\propto(1+z)$ are calculated in the metric language. A finite-size imprint at the horizon scale should appear in low-multipole power and in cosmological topology.

\subsection{Discreteness and the \texorpdfstring{$C_3$}{C3} Symmetry of the Generations}

In the Standard Model, the masses of the electron, muon, and tau lepton are empirical parameters. RefG reads them as an internal, order-nine, three-phase ($C_3$) triplet of a single topological defect: three distinct resonant modes of one physical structure.

The parametrization of the charged leptons combines the relation $m\propto\nu^2$, or $\nu\propto\sqrt{m}$, with a three-phase structure:

\begin{equation}
\nu_k=\mu\left[1+A_K\cos\left(\theta+\frac{2\pi k}{3}\right)\right],
\qquad k=0,1,2.
\label{eq:c3-triplet}
\end{equation}

Here $\mu$ is the common frequency scale, $A_K$ is the amplitude of the three-phase modulation, $\theta$ is the initial phase, and $k$ labels the three generations. In RefG's candidate order-nine return map, $\theta=h/9$; the selected $h=2$ branch gives $\theta=2/9$ radians. The relation $m\propto\nu^2$ reads mass as the integrated energy of a locally confined oscillation.

For $A_K=\sqrt{2}$, the $C_3$ algebra gives the empirical Koide ratio, $K=2/3$, exactly \citep{koide1982,koide1983}. Both the three-phase parametrization and its geometric reading are known in the literature: Brannen recast it in the language of a circulant phase \citep{brannen2006}, while Foot related $K=2/3$ to the $45^\circ$ geometry of the $\sqrt{m}$ vector \citep{foot1994}. RefG further reads $2/9$ as the reduced phase of the defect's order-nine return cycle.

The value $A_K=\sqrt{2}$ is associated with equalization of the internal channels, while $\theta=2/9$ corresponds to the $h=2$ branch of the order-nine cycle. Anchored to the electron pole mass, the exact branch $(A_K,\theta)=(\sqrt{2},2/9)$ misses the muon mass by approximately $451.7\sigma$ and the tau mass by $0.61\sigma$. This discrepancy shows directly that the radiative protection of the geometric picture, and the bridge from that picture to pole masses, require a separate quantitative explanation \citep{sumino2009}.

This balance also has a simple musical intuition: the charged-lepton $C_3$/Koide triplet resembles the internal spectrum of a three-dimensional volumetric resonator more than it resembles an ordinary harmonic chord. The tones here are different modes of one body.

A complete enumeration of all stable charged modes will decide whether a fourth charged lepton exists.

\section{Topology and Quantum Numbers}

\subsection{Phase, Orientation, and Topological Quantization}

Up to this point, we have spoken of the universal underlying medium primarily in the language of amplitude and local pressure. Its full local state is described by a collection of fields. The field $\Phi$ is a phase--clock scalar field; separate variables describe the amplitude profile and the spatial frame. Together, they carry amplitude, phase, internal form, and orientation.

If $\Phi$ is a single-valued scalar and the flow is defined directly by $u=\nabla\Phi$, the vorticity vanishes. A topological sector requires a multivalued phase, a defect line, and an internal orientation locked in space. The amplitude vanishes at the defect core, while the phase winds around it an integer number of times. A quantized vortex in a superfluid is a laboratory example of just such a structure.

The cosmological counterpart of this structure is the \textbf{cosmic string}---a macroscopic defect line frozen in during a phase transition in the early Universe. Its formation is described by Kibble's causal mechanism \citep{kibble1976}. In the language of RefG, this is the ``frozen seam'' mentioned in Chapter IV. The vacuum topology of the underlying medium determines whether such lines form, while observations constrain their cosmological density.

Lord Kelvin's work on vortex atoms and vortex motion \citep{thomson1867,thomson1869} led Peter Tait to catalog aether knots and helped give rise to the mathematical theory of knots \citep{tait1877}. RefG carries this historical line forward through pressure--energy balance and population-wide resonant locking.

Phase quantization arises when the field is continuous along a loop, the state returns to a single-valued configuration after one complete circuit, and the loop belongs to a nontrivial topological class. Under these conditions, the total change in phase is an integer multiple of $2\pi$.

The Călugăreanu--White--Fuller relation applies to a closed, orientable, framed ribbon \citep{calugareanu1961,white1969,fuller1971,ricca2011}:

\begin{equation}
\operatorname{Lk}=\operatorname{Wr}+\operatorname{Tw}
\label{eq:self-linking}
\end{equation}

Here $\operatorname{Lk}$ is the integer linking number, $\operatorname{Wr}$ describes the spatial writhe of the axis, and $\operatorname{Tw}$ describes the internal twist of the frame. The geometrical contributions $\operatorname{Wr}$ and $\operatorname{Tw}$ may change, but their sum $\operatorname{Lk}$ remains invariant under the corresponding continuous deformations.

In hydrodynamics, helicity is a conserved quantity in an ideal fluid, and Moffatt related it to the linking number \citep{moffatt1969}. Onsager and Feynman predicted the quantization of circulation in superfluid helium, and Vinen measured the quantum of circulation experimentally \citep{onsager1949,feynman1955,vinen1961}. Integer phase winding is thus a directly established laboratory phenomenon.

In the RefG map, the formula \textbf{``topological locking $=$ charge''} means that electric charge is the external readout of a locked phase--orientation path. The topological number fixes the discrete class, while the field dynamics describe its electromagnetic readout.

The Möbius strip does not lie within the direct domain of the Călugăreanu--White--Fuller relation; it is described through an orientable double cover. We shall use the spinorial intuition of this double cover in the next subsection.

This direction also offers a testable indication. Kelvin waves have been directly excited on quantized vortices in superfluid helium, and their dispersion has been measured \citep{minowa2025}. RefG expects the internal Kelvin waves of an oscillon to have a corresponding topological dispersion: the radial mode is linked to a Compton-like mass pulsation, and the helical mode to the radiation channel.

\subsection{Spin and Fermionic Statistics}

A half-turn twist of the phase along an oscillon's internal axis creates a double-cover structure like a Möbius strip. One circuit reverses the orientation; two restore the initial state. This $Z_2$ double cover bears the defining signature of a spinorial transformation: a $2\pi$ rotation changes the sign of the state, while a $4\pi$ rotation restores it completely. RefG associates the electron's spin-$1/2$ with this geometry. Changing the topological class requires the amplitude to vanish at the core, the defect to reconnect, or the boundary condition to change.

This geometry also gives an intuition for fermionic exchange. Exchanging two identical double-cover modes must multiply the joint state by minus one, making the two-particle wave function antisymmetric. Placing both particles in the same one-particle state then gives a vanishing two-particle state---a geometrical picture of the Pauli principle. An initial version of this vortex-exchange route was developed in the author's earlier work \citep{lotishvili2026b}.

In a bosonic mode, the exchange phase does not change the sign of the state, so many bosons can occupy the same quantum state. The phase-locked, or coherent, photonic state of a laser is a vivid example.

This remains a geometrical intuition. In relativistic quantum field theory, the spin--statistics relation is a rigorous theorem based on Lorentz invariance, positivity of energy, and microcausality \citep{pauli1940}. A complete RefG derivation must recover these conditions, together with the antisymmetric construction of Fock space.

\subsection{Charge, the Photon, and the Electromagnetic Channel}

RefG reads electric charge as the external imprint of an oscillon's internal phase twist. The two opposite orientations of the twist correspond to positive and negative charge. Outside the oscillon, this imprint spreads in every direction, and its total flux through a closed surface determines the magnitude of the charge.

In this map, the positron is the topological mirror of the electron. When a particle and its antiparticle annihilate, their opposite phase windings combine, and the released energy passes into the radiation channel. In a neutral mode, the contributions of opposite phase twists balance, and the net flux through a closed surface vanishes. RefG associates the neutrino with this neutral sector.

Turning the asymmetry angle by the same amount at every point leaves the physical picture unchanged. A measurable imprint arises when this angle varies in space. RefG reads the photon as a transverse, helical wave of that variation: it carries changes in the asymmetry phase and mediates the electromagnetic interaction.

Topology also appears in classical Maxwell theory. The field lines of Rañada's vacuum solutions form linked loops and are classified by the Hopf index \citep{ranada1989}. RefG connects this mathematical picture to the underlying-medium channel of charge and electromagnetic dynamics.

The sign of the force is determined jointly by the mediator's spin, the type of source, and the sign of the coupling. In the electromagnetic vector channel, like charges repel and opposite charges attract; in the tensor gravitational channel, positive energy produces universal attraction. In three spatial dimensions, the dilution of a spherical flux gives a $1/r^2$ law, while the corresponding potential follows a $1/r$ dependence.

In this system, the fine-structure constant $\alpha$ is the electromagnetic readout of the phase--topological locking of a charged oscillon. From this same channel, the quantitative sector must recover Maxwell's equations, Gauss's law, the two photon helicities, and charge conservation.

\subsection{Color, Confinement, and Three Generations}

Three-dimensionality plays a structural role in the RefG particle map. The three spatial directions serve as a geometrical carrier for the internal color triplet. Color remains an internal degree of freedom; the spatial axes provide its geometrical basis.

A structure oscillating along a single axis is topologically incomplete and has open ends. Combining all three axes corresponds to baryonic closure; joining an axis to an anti-axis corresponds to mesonic closure. An isolated quark cannot form a closed state, and RefG associates confinement with this geometry. Quarks carry three color states, gluons mediate transitions between them, and leptonic modes are color-neutral.

Here $D_3$ denotes the discrete deformation used in the eight-mode count. The resulting eight modes numerically match the eight generators of $SU(3)$ and form a structural map of the color sector. Full $SU(3)$ requires these modes to reproduce its algebra, local symmetry, and couplings.

The question of three generations is the logical continuation of the $C_3$ triplet discussed in the preceding chapter. In the RefG reading, charged leptons are the internal three-phase, order-nine modes of a single topological defect, while the quark generations carry the same generational architecture into the colored axial sector. At this level, a generation is no longer a replication added to the table from outside; it is a recurring mode of internal topological locking in the underlying medium.

In RefG, three-dimensionality is tied to the stability of knotted loops. One-dimensional closed loops form nontrivial knots precisely in three spatial dimensions; higher dimensions provide additional directions through which they can be untied. This argument singles out three-dimensionality in models where matter is built from such loops \citep{ehrenfest1918,whitrow1955,berera2017}.

\subsection{A Geometrical Map of the Standard Model}

\subsubsection{The Electron and Leptons}

In this map, the electron is a complete toroidal structure with a Möbius-type twist, in which all three spatial axes close locally. RefG associates this double-cover geometry with the Dirac value $g=2$. The measured value of $g_e$ also contains radiative corrections from quantum electrodynamics \citep{schwinger1948,fan2023}; a topological model must reproduce those corrections quantitatively.

\subsubsection{Quarks and Color Charge}

In this map, quarks are axial segments with open ends. Their inability to exist in isolation corresponds to confinement, while fractional charge is read as the topological phase locked in thirds. Color is an internal axial asymmetry, and the gluon carries transitions between color states.

\subsubsection{Photons and Bosons}

The photon is a massless, helical transverse wave without localized confinement. RefG associates its two helicities and $U(1)$ structure with the transverse mode of the underlying medium. The eight-mode map of the gluons was described in the preceding subsection. In the RefG map, the $W$ and $Z$ bosons are locally confined, massive mediator modes.

\subsubsection{Neutrinos and the Higgs Boson}

In this map, the neutrino is an electrically neutral, incompletely confined mode. RefG reads the Higgs boson as a nonrotating, spherically symmetric radial pulsation, naturally corresponding to a spin-0 state. The quantitative electroweak sector must recover neutrino masses and mixing, the $125\,\mathrm{GeV}$ Higgs mass, the vacuum expectation value, and the Yukawa couplings.

\subsubsection{The Graviton and the Weinberg--Witten Theorem}

The local tensor branch of RefG contains two polarizations, $h_+$ and $h_\times$. The Weinberg--Witten theorem places strict constraints on a Lorentz-covariant conserved tensor and on composite massless high-spin states \citep{weinbergwitten1980}. RefG's fundamental medium-based description examines the theorem's assumptions separately; in the low-energy limit, a massless tensor mode, universal coupling, and effective Lorentz symmetry must emerge.

\subsubsection{Dynamics and the Lagrangian}

The mathematical language of this map is the Lagrangian---the dynamical law from which the field equations and admissible forms follow. Even a simple local field model can admit stable knotted solitons \citep{faddeev1997}. RefG connects this mechanism to the particle spectrum, masses, and decay rules.

\subsection{Fundamental Constants: \texorpdfstring{$\alpha$}{alpha} and the Gravitational Hierarchy}

The fine-structure constant, $\alpha\approx1/137$, is one of the celebrated mysteries of physics---Feynman regarded it as one of the great unexplained numbers of theoretical physics \citep{feynman1985}. RefG reads $\alpha$ as the relation between the internal phase--topological locking of a charged oscillon and its external electromagnetic readout.

The effective boundary calculation gives $\alpha^{-1}=137.035999174\ldots$ \citep{lotishvili2026a}. This is a target-fitted closure: the observed value of $\alpha$ and the lepton mass ratios were used to select the branch and the law. Its significance must be tested by a new, independent prediction.

The second fundamental question is the \textbf{hierarchy problem}: why is gravity so weak? Between two protons, the electromagnetic force is roughly $10^{36}$ times stronger than gravity; for a pair of electrons, the ratio is of order $10^{42}$. The relation between these large numbers became the historical starting point of Dirac's hypothesis \citep{dirac1937}.

RefG traces the hierarchy to the structural difference between two channels: \textbf{direct phase--transverse coupling} and a \textbf{secondary pressure imprint}.

\begin{itemize}
\item The \textbf{Planck mass}, $M_P=\sqrt{\hbar c/G}$, is the calibration scale for gravitational coupling.
\item The \textbf{electron and proton} are microscopic oscillons with masses many orders of magnitude below the Planck mass.
\end{itemize}

Here gravity is the long-range imprint of the pressure deficit created by an oscillon's energy, whereas the electromagnetic interaction is a direct phase--transverse channel.

For two electrons, the force ratio can be written algebraically as

\begin{equation}
\frac{F_{\mathrm{em}}}{F_g}=\alpha\left(\frac{M_P}{m_e}\right)^2.
\label{eq:force-hierarchy}
\end{equation}

This formula displays the hierarchy through the ratio of the Planck mass to the electron mass; the common origin of $G$, $M_P$, and the microscopic energy of the oscillon requires a separate dynamical explanation.

In this map, the weakness of gravity reflects the different ontological nature of the two interaction channels. A quantitative closure must unite particle masses, the photon mode, and charge normalization in a single system.

\section{Epilogue: Limits of Knowledge}

\subsection{Why a Single Underlying Medium?}

Let us return, in closing, to the question: why a single underlying medium? One motivation is the principle of parsimony---a common carrier reduces the number of independent starting assumptions. Yet simplicity alone cannot select a physical theory.

The more important argument concerns the \textbf{origin of regularities}. By treating the electron and the photon as different resonant modes of the same underlying medium, RefG seeks to derive the conservation of energy and charge from the dynamical balance of their common carrier, rather than leave these laws as independent starting assumptions. In this language, interactions are couplings between modes of the common medium, and conservation laws are geometrical expressions of their balance.

This argument has a historical background. Eugene Wigner wrote of the ``unreasonable effectiveness'' of mathematics in the natural sciences \citep{wigner1960}. The ontology of a common carrier offers a possible physical explanation for this order: shared mathematical regularities may reflect a shared structure. Deriving this connection dynamically is itself a task of the research program.

The argument also has an empirical side: the problem of precise mutual compatibility. The electrical neutrality of macroscopic matter, the stringent upper bound on the photon mass, and the pattern of fractional quark charges call for a common structural explanation \citep{navas2024}. RefG connects this order to the geometry of modes of a single underlying medium.

A single underlying medium gives rise to a rich phenomenology: it can carry pressure, shear, phase, topology, rotation, memory, and resonance at once. This position echoes Philip Anderson's well-known dictum, ``More is different'': every new level of organization has regularities of its own \citep{anderson1972}. In RefG, the transitions between levels must be calculated from this multichannel underlying medium.

One of RefG's cosmological motivations is the picture of a finite, boundaryless global resonator. The discrete spectrum discussed in Chapter IV points in a direction compatible with this hypothesis. Its physical content is determined by cosmological solutions and observational signatures.

\subsection{Three Operational Infinities}

The Universe is always measured by something that is itself part of the Universe: our clocks, rulers, detectors, and bodies are all oscillon modes of the universal underlying medium. This idea of a ``self-measuring Universe'' is related to Wheeler's ``It from Bit'' program \citep{wheeler1990}. For Wheeler, information is primary; in RefG, the underlying medium as substance is primary, and the informational imprint is a result of its self-differentiation.

In what follows, ``infinity'' has an operational meaning: it denotes a limit or an unbounded extension of a measurement procedure, without presupposing that the Universe is literally infinite in size, age, or divisibility.

Here Gödel's theorem supplies only a logical analogy; it establishes no physical conclusion. In any sufficiently strong, consistent formal system, there are propositions that cannot be decided within that system \citep{godel1931}. A similar question arises in a self-measuring Universe: how completely can one part of the system measure the boundaries of the whole? From this intuition emerge the ``three operational infinities.''

\textbf{Temporal infinity:} Before the first resonant clocks, operational time was undefined. The figure ``13.8 billion years'' is metric time calibrated by present-day clocks. A process-time interpretation may assign a longer duration to the internal history of the early underlying medium; this is a limit of what clocks can measure, not a claim of an additional cosmological time.

\textbf{Spatial infinity:} We do not measure ``pure'' emptiness; we measure only the connections between distinguishable nodes. The cosmological horizon is therefore the operational boundary of our measurable network, not the physical end of the underlying medium itself. The region beyond the horizon remains outside direct observation.

\textbf{Scale infinity:} In RefG, the Planck length denotes not the absolute floor of the underlying medium, but a natural crossover scale built from $c$, $\hbar$, and $G$, where quantum and gravitational effects become important together. In the history of physics, the ``smallest'' known scale has changed more than once---from the atom to the nucleon, and from the nucleon to the quark. The physical content of still smaller scales belongs to a further dynamical layer of the underlying medium.

In all three cases, infinity arises between the measurement procedure and the demand for a complete description. Even a finite object can unfold without bound in a chosen representation: stereographic projection maps a finite sphere onto an infinite plane. Mandelbrot's ``coastline paradox'' expresses a related idea, with the measured length increasing as the measuring scale decreases \citep{mandelbrot1982}. RefG reads these examples as operational properties of description and measurement scale.

\subsection{Methodological Principle}

The methodological principle of this work is clear: an intuitive physical picture must become a rigorous, calculable, and falsifiable formalism. The analysis of a particle oscillon follows this sequence:

\begin{enumerate}
\item \textbf{Experimental basis:} record the mass, charge, spin, lifetime, and other observables.
\item \textbf{Mode identification:} select the corresponding resonant class of the underlying medium.
\item \textbf{Initial geometrical form (ansatz):} describe the oscillon's form, topology, and boundary conditions precisely.
\item \textbf{Dynamical law:} derive the action, field equations, and source.
\item \textbf{Symmetries and conservation laws:} determine the charge, spin, color, and zero modes.
\item \textbf{Existence and stability:} test the finite-energy solution, matching, and perturbation spectrum.
\item \textbf{Parameter closure:} determine every coefficient independently of the data we seek to predict.
\item \textbf{New prediction:} calculate an independent observable and its uncertainty.
\item \textbf{Data test:} compare the result with observations using a predefined falsification criterion.
\end{enumerate}

The aim of this sequence is to derive the gyromagnetic factor from oscillon dynamics, and to turn the formula ``topological locking $=$ charge'' into a measurable, independently testable law.

\section*{Conclusion}
\addcontentsline{toc}{section}{Conclusion}

This article has built a single, continuous conceptual chain grounded in the principles of Refractive Gravity (RefG). Its ontological picture of the Universe is this: the spontaneous self-differentiation, or self-distinction, of one underlying medium gives rise to measurability---the operational seeds of space, time, and mass. A locally confined resonant form of the medium, the \textbf{oscillon}, creates a pressure deficit around itself that an external observer reads as effective mass. Motion within the effective geometry altered by this deficit is gravity: not a mechanical force, but the universal refractive response to the state of the medium.

This unified pressure chain unfolds across three principal scales:

\begin{itemize}
\item at the \textbf{local scale}---in compact objects, where internal state and external readout are clearly separated;
\item at the \textbf{galactic scale}---in the large-scale vortical flow of the underlying medium, which RefG connects to MOND phenomenology;
\item at the \textbf{cosmological scale}---in the historical, asymptotic relaxation of the state of the underlying medium.
\end{itemize}

The allowed spectrum of the underlying medium and mutual locking across the population connect the sparse, discrete particle spectrum to a single mechanism, while the topology of phase and orientation gives quantum numbers a geometrical language. Along the same line, the effective boundary readout of the fine-structure constant $\alpha$ connects the oscillon map to the numerical structure of the electromagnetic interaction.

The article rests on the weak-field regime---the domain in which gravity has been tested most precisely. On the selected 1PN branch, the refractive language of RefG agrees with the standard description of general relativity. From this foundation, the research program extends to particle, galactic, and cosmological scales, where the physical map brings completed calculations and candidate observational tests into one system.

From a single ontological foundation, Refractive Gravity develops a research program whose aim is a unified, calculable, and empirically testable description of gravity, matter, and measurability. The present text lays the foundation of this physical paradigm and opens the way for further research.

\section*{Current Status of the Research Program}
\addcontentsline{toc}{section}{Current Status of the Research Program}

The statuses are collected here in one place so that the main text can preserve the flow of the physical narrative.

\begin{enumerate}
\item \textbf{The underlying medium and the origin of measurability.} The ontological core has been formulated; operational distinguishability defines the starting point from which space, time, and mass emerge.

\item \textbf{Inter-node separation and minimum distinguishability.} The idea of operational distinguishability has been formulated; a threshold of the order of the Planck scale is a physical candidate and has not been derived from the action.

\item \textbf{The local collective rate.} Common-transposition invariance and toy models of mutual locking have been obtained; calculation of the actual stable set of oscillons remains open.

\item \textbf{Oscillon deficit, size, time, and mass.} The metric filters and the first-order biconformal branch have been obtained symbolically; a finite-energy nonlinear oscillon profile and the complete mass--readout map remain open.

\item \textbf{Weak-field gravity.} On the selected 1PN branch, $\gamma=1$, $\beta=1$, and the listed observables agree with the standard results; higher orders and a complete likelihood-based test against the data remain open.

\item \textbf{Inertia and gravitational mass.} The Noether identity has been obtained conditionally for a finite-energy, Lorentz-covariant localized solution satisfying the Laue boundary condition; a direct derivation of the zero mode remains open.

\item \textbf{Gravitational waves and local unobservability.} In the local Minkowski limit, the tensor propagation speed and the transverse--traceless detector response have been obtained; cosmological propagation, source--flux normalization, the complete waveform, and scalar polarization remain open.

\item \textbf{Compact objects.} A static, conditional map has been obtained; the domains of the exterior and core solutions do not yet overlap. A finite core, collapse dynamics, an equation of state, rotation, and full stability remain open.

\item \textbf{The galactic vortex and MOND/Tully--Fisher.} The transition algebra and the baryonic Tully--Fisher relation (BTFR) follow conditionally once $a_0=cH/(2\pi)$ is assumed and vortex closure is imposed; derivation from the action, SPARC/RAR tests, and complete modeling of clusters remain open.

\item \textbf{Cosmological relaxation and dark energy.} The Friedmann--Lemaître--Robertson--Walker (FLRW) stress algebra and a conditional $\Lambda$-like diagnostic have been obtained; pressure relaxation is an interpretive bridge. Perturbations and numerical fits to BBN, CMB, Planck, and BAO data remain open.

\item \textbf{Two-level time and quantum randomness.} The interpretive map is compatible with the facts of Bell tests and delayed-choice experiments; the Born distribution, the no-signaling condition, and the complete dynamics of measurement have not been derived.

\item \textbf{Particles as oscillon modes.} The geometrical placement map and initial spectral scheme have been formulated as candidates; a finite-energy oscillon, its form factor, masses, couplings, and decays remain open.

\item \textbf{The Koide/$C_3$/$2/9$ structure.} The $C_3$ algebra gives $K=2/3$ exactly; with $(A_K,\theta)=(\sqrt{2},2/9)$ and the electron pole mass as the anchor, the result misses the muon mass by $451.7\sigma$ and the tau mass by $0.61\sigma$. The choice $h=2$, the relation $m\propto\nu^2$, derivation from the action, and radiative protection remain open.

\item \textbf{Spin, Pauli exclusion, and fermionic statistics.} The Möbius double-cover geometry provides a candidate mechanism; a complete derivation of Fock space and of the spin--statistics relation of quantum field theory remains open.

\item \textbf{Charge, the photon, and the electromagnetic channel.} Subtheorems establishing local-frame $U(1)$ symmetry, two helicities, a screened far field, and a conserved phase current have been obtained; canonical charge normalization, Maxwell's equations, Gauss's law, the Ward identities, quantum electrodynamics (QED), and the Dirac sector remain open.

\item \textbf{Color, quarks, and gluons.} The three-color, eight-mode construction provides only a partially compatible map; the structure constants of $SU(3)$, interaction vertices, scale evolution, the confinement potential, and fractional charge remain open.

\item \textbf{The electroweak map and the Higgs.} This is a geometrical placement map; $SU(2)_L$ chirality, a $W/Z$ mass matrix that preserves a massless photon, CKM/PMNS mixing, the neutrino mass ordering, Yukawa couplings, and the vacuum expectation value have not been derived.

\item \textbf{The fine-structure constant $\alpha$.} A target-fitted closure of the effective boundary action has been obtained. The final formula uses $h=2$, the lepton ratios, and $N=34$; the observed value of $\alpha$ was used to select the branch and the law. The microscopic core, the internal leptonic threshold, and an independent observable remain open.

\item \textbf{The gravitational hierarchy.} The force ratio can be written algebraically as $F_{\mathrm{em}}/F_g=\alpha(M_P/m)^2$; the origins of $G$, $M_P$, the oscillon's microscopic energy scale, and the mass hierarchy remain open.

\item \textbf{A finite self-measuring Universe.} The global resonator is a cosmological hypothesis; low-multipole and topological signatures are candidates, and no unique likelihood fit to the data has been obtained.
\end{enumerate}

\phantomsection
\addcontentsline{toc}{section}{References}
\bibliographystyle{plainnat}
\bibliography{Intuitive_Article_EN}

\end{document}
